\chapter{Ordens de grandeza: medir o universo}\label{ch:g10:orders-of-magnitude}

A física começa onde termina a opinião: em uma medida. Mas as coisas
que a física mede cobrem um intervalo absurdo --- um núcleo atômico
mede cerca de \qty{e-15}{m} de ponta a ponta e o universo observável
cerca de \qty{e27}{m}: um fator $10^{42}$, muito além do que qualquer
imagem mental consegue abrigar. Este capítulo instala as ferramentas
que domam esse intervalo --- \omterm{def:g10:orders-of-magnitude:unit}{unidades}, potências de dez, \omterm{def:g10:orders-of-magnitude:sigfig}{algarismos
significativos} --- e dois hábitos usados em todas as páginas
seguintes: escrever a \omterm{def:g10:orders-of-magnitude:unit}{unidade} e verificar a \omterm{def:g10:orders-of-magnitude:oom}{ordem de grandeza}.

\section{As unidades: a gramática da medida}

\begin{definition}[Grandezas físicas e unidades do SI]\label{def:g10:orders-of-magnitude:unit}
Uma \emph{grandeza física}\index{grandeza física} é uma propriedade do
mundo que se pode medir: um comprimento, uma duração, uma massa. Medi-la
é contar quantas vezes cabe nela uma grandeza de referência --- a
\emph{unidade}\index{unidade}. O \emph{Sistema Internacional de
Unidades}\index{Sistema Internacional de Unidades (SI)} (SI) fixa sete
unidades de base, das quais três governam este volume inteiro: o
\emph{metro}\index{metro} (\unit{m}) para o comprimento, o
\emph{quilograma}\index{quilograma} (\unit{kg}) para a massa e o
\emph{segundo}\index{segundo} (\unit{s}) para o tempo. (As outras quatro
--- ampère, kelvin, mol, candela --- entram em cena quando chegarem a
eletricidade, o calor e a química.)
\end{definition}

\begin{remark}[Um número sem unidade não é um resultado]\label{rem:g10:orders-of-magnitude:units}
``A distância é $384$'' não diz nada: \qty{384}{mm}, \qty{384}{km} e
\qty{384}{Mm} são três mundos diferentes (o último é, por acaso, a
distância até a Lua). Em 1999 a NASA perdeu o \emph{Mars Climate
Orbiter}, de $327$ milhões de dólares, porque um programa calculava as
forças dos propulsores em libras-força e outro as lia como newtons.
Regra da casa, desta linha até o fim do livro: todo número medido ou
calculado carrega a sua \omterm{def:g10:orders-of-magnitude:unit}{unidade}, em cada linha do cálculo.
\end{remark}

\begin{definition}[Prefixos do SI]\label{def:g10:orders-of-magnitude:prefixes}
Um \emph{prefixo}\index{prefixo (SI)} colado a uma \omterm{def:g10:orders-of-magnitude:unit}{unidade} multiplica-a
por uma potência de dez fixa:
\begin{center}
\begin{tabular}{llc@{\qquad\qquad}llc}
femto & f & $10^{-15}$ & centi & c & $10^{-2}$ \\
pico  & p & $10^{-12}$ & quilo & k & $10^{3}$  \\
nano  & n & $10^{-9}$  & mega  & M & $10^{6}$  \\
micro & $\mu$ & $10^{-6}$ & giga & G & $10^{9}$ \\
mili  & m & $10^{-3}$  & tera  & T & $10^{12}$ \\
\end{tabular}
\end{center}
Assim, $\qty{1}{nm} = \qty{e-9}{m}$ e $\qty{1}{km} = \qty{e3}{m}$.
\end{definition}

\begin{example}[Ler os prefixos]\label{ex:g10:orders-of-magnitude:prefixes}
A luz verde tem comprimento de onda de cerca de
$\qty{500}{nm} = \qty{5.00e-7}{m}$; um fio de cabelo humano tem cerca de
$\qty{70}{\micro m} = \qty{7.0e-5}{m}$ de espessura; a Lua orbita a
$\qty{384}{Mm} = \qty{3.84e8}{m}$. Uma gota de chuva tem massa de cerca
de $\qty{30}{mg} = \qty{3.0e-5}{kg}$ --- repare que a \omterm{def:g10:orders-of-magnitude:unit}{unidade} de base do
SI para a massa é a única das sete que já vem com \omterm{def:g10:orders-of-magnitude:prefixes}{prefixo}.
\end{example}

\section{Notação científica e algarismos significativos}

\begin{definition}[Notação científica]\label{def:g10:orders-of-magnitude:scinot}
Um número está em \emph{notação científica}\index{notação científica}
quando é escrito na forma
\[
a \times 10^{n}, \qquad 1 \leq a < 10, \quad n \text{ inteiro.}
\]
O fator $a$ é a \emph{mantissa}\index{mantissa} e o inteiro $n$ é o
\emph{expoente}\index{expoente}. Assim,
$\num{149600000} = \num{1.496e8}$ e $\num{0.00052} = \num{5.2e-4}$.
\end{definition}

\begin{proposition}[Calcular com potências de dez]\label{prop:g10:orders-of-magnitude:powers}
Para quaisquer inteiros $m$ e $n$:
\[
10^{m} \times 10^{n} = 10^{m+n}, \qquad
\frac{1}{10^{n}} = 10^{-n}, \qquad
\left(10^{m}\right)^{n} = 10^{mn}.
\]
\end{proposition}

\begin{proof}[Demonstração parcial]
Para \omterm{def:g10:orders-of-magnitude:scinot}{expoentes} positivos, basta contar os fatores: $10^{m} \times 10^{n}$
é um produto de $m + n$ dezes, e $(10^m)^n$ encadeia $n$ grupos de $m$
dezes. Escrever $10^{-n}$ no lugar de $1/10^{n}$ é justamente a convenção
que mantém a primeira regra válida para todos os inteiros: ela obriga
$10^{n} \times 10^{-n} = 10^{0} = 1$. O volume de matemática
correspondente estuda as regras dos \omterm{def:g10:orders-of-magnitude:scinot}{expoentes} por completo.
\end{proof}

\begin{example}[Multiplicar em notação científica]\label{ex:g10:orders-of-magnitude:scinot}
Quantas Terras cabem no Sol, em massa? Com
$m_{\text{Sol}} = \qty{1.99e30}{kg}$ e
$m_{\text{Terra}} = \qty{5.97e24}{kg}$:
\[
\frac{\num{1.99e30}}{\num{5.97e24}}
= \frac{1.99}{5.97} \times 10^{30-24}
= 0.333 \times 10^{6}
= \num{3.33e5}.
\]
\omterm{def:g10:orders-of-magnitude:scinot}{Mantissas} e \omterm{def:g10:orders-of-magnitude:scinot}{expoentes} são tratados separadamente; a resposta final é
renormalizada para que a \omterm{def:g10:orders-of-magnitude:scinot}{mantissa} volte a cair em $[1, 10)$.
\end{example}

\begin{definition}[Algarismos significativos]\label{def:g10:orders-of-magnitude:sigfig}
Os \emph{algarismos significativos}\index{algarismos significativos} de
um número medido são os dígitos por que a medida realmente responde:
todos os dígitos escritos, menos os zeros à esquerda. Os zeros à direita
contam --- $\qty{5.20e-3}{s}$ (três algarismos significativos) promete
mais do que $\qty{5.2e-3}{s}$ (dois). A \omterm{def:g10:orders-of-magnitude:scinot}{notação científica} mostra a
contagem de relance: a \omterm{def:g10:orders-of-magnitude:scinot}{mantissa} carrega exatamente os algarismos
significativos.
\end{definition}

\begin{method}[Quantos dígitos guardar]\label{met:g10:orders-of-magnitude:sigfig}
Um resultado calculado não pode ser mais preciso do que os seus
ingredientes.
\begin{enumerate}
  \item Em um produto ou em um quociente, guarde tantos \omterm{def:g10:orders-of-magnitude:sigfig}{algarismos
        significativos} quantos tiver o fator \emph{menos} preciso.
  \item Arredonde apenas a resposta final, nunca as etapas
        intermediárias.
  \item Na dúvida, três \omterm{def:g10:orders-of-magnitude:sigfig}{algarismos significativos} são a precisão de
        trabalho deste livro.
\end{enumerate}
\end{method}

\begin{example}[A circunferência da Terra]\label{ex:g10:orders-of-magnitude:sigfig}
O diâmetro da Terra é $d = \qty{1.27e7}{m}$ (três \omterm{def:g10:orders-of-magnitude:sigfig}{algarismos
significativos}), de modo que a sua circunferência vale
$\pi d = \num{3.14159} \times \qty{1.27e7}{m} = \qty{3.99e7}{m}$ ---
três algarismos, e não os oito que a calculadora exibe. Escrever
$\qty{39898227}{m}$ seria alegar conhecer o planeta com precisão de um
\omterm{def:g10:orders-of-magnitude:unit}{metro}.
\end{example}

\begin{omfigure}
\begin{tikzpicture}[x=1.9cm, font=\small]
  \draw[-Stealth] (-0.3,0) -- (5.45,0);
  \foreach \e in {0,...,5} {
    \draw (\e,-0.1) -- (\e,0.1);
    \node[below] at (\e,-0.12) {$10^{\e}$};
  }
  \foreach \x in {2.301,2.477,2.602,2.699,2.778,2.845,2.903,2.954}
    \draw (\x,-0.05) -- (\x,0.05);
  \draw[Stealth-Stealth, omMeth] (0,0.45) -- (1,0.45)
    node[midway, above] {$\times 10$};
  \fill[omDef] (2.477,0) circle (1.7pt);
  \node[above=3pt, omDef] at (2.477,0.06) {$3 \times 10^{2}$};
  \draw[-Stealth, omDef, dashed] (2.38,0.22) -- (2.06,0.22);
  \fill[omThm] (3.845,0) circle (1.7pt);
  \node[above=3pt, omThm] at (3.845,0.06) {$7 \times 10^{3}$};
  \draw[-Stealth, omThm, dashed] (3.90,0.22) -- (3.96,0.22);
\end{tikzpicture}

{\small Uma régua em que cada passo igual \emph{multiplica} por dez. Os
traços pequenos são $200, 300, \dots, 900$: uma década inteira se
espreme em cada segmento. Nessa régua $\num{300}$ pende para $10^{2}$ e
$\num{7000}$ pende para $10^{4}$ --- as suas \omterm{def:g10:orders-of-magnitude:oom}{ordens de grandeza}.}
\end{omfigure}

\section{A escada do universo}

\begin{definition}[Ordem de grandeza]\label{def:g10:orders-of-magnitude:oom}
A \emph{ordem de grandeza}\index{ordem de grandeza} de uma grandeza é a
potência de dez mais próxima dela: escreva a grandeza como
$a \times 10^{n}$ com $1 \leq a < 10$ e tome $10^{n}$ se $a < 5$ e
$10^{n+1}$ se $a \geq 5$. (A convenção exata da fronteira nunca importa:
uma afirmação de ordem de grandeza só vale mesmo a menos de um fator
dez.)
\end{definition}

\begin{example}[Ordens de grandeza]\label{ex:g10:orders-of-magnitude:oom}
Uma vida humana dura cerca de $80$ anos, ou seja,
$80 \times \qty{3.16e7}{s} \approx \qty{2.5e9}{s}$: \omterm{def:g10:orders-of-magnitude:oom}{ordem de grandeza}
$\qty{e9}{s}$ --- alguns bilhões de segundos, gaste-os bem. Um ser
humano tem $\qty{1.7}{m}$ de altura: ordem $\qty{e0}{m}$. O monte
Everest, $\qty{8848}{m}$: ordem $\qty{e4}{m}$, já que $8.848 \geq 5$.
\end{example}

Empilhar essas estimativas por \omterm{def:g10:orders-of-magnitude:scinot}{expoente} produz a imagem mestra deste
capítulo: uma escada em que cada degrau é um mundo dez vezes mais largo
do que o de baixo.

\begin{omfigure}
\begin{tikzpicture}[x=0.26cm, font=\small]
  \draw[-Stealth] (-16.5,0) -- (28.7,0);
  \foreach \e in {-15,-10,-5,0,5,10,15,20,25} {
    \draw (\e,-0.12) -- (\e,0.12);
    \node[below, font=\footnotesize] at (\e,-0.14) {$10^{\e}$};
  }
  \foreach \e/\t in {-15/núcleo, -10/átomo, -6/bactéria, 0/humano,
                     7/Terra, 11/Terra--Sol, 16/ano-luz,
                     21/Via Láctea, 27/universo} {
    \draw[omDef, thick] (\e,0.14) -- (\e,0.62);
    \node[rotate=52, anchor=south west, inner sep=1pt, omDef]
      at (\e,0.66) {\t};
  }
  \node[below, font=\footnotesize] at (26.6,-0.75) {tamanho (\unit{m})};
\end{tikzpicture}

{\small A escada do universo, em \omterm{def:g10:orders-of-magnitude:unit}{metros}: $42$ potências de dez separam o
núcleo atômico do universo observável, e o ser humano está a dez degraus
do átomo, mas a vinte e sete do topo.}
\end{omfigure}

Acima do degrau da Terra os \omterm{def:g10:orders-of-magnitude:unit}{metros} deixam de ser cômodos; a astronomia
tem \omterm{def:g10:orders-of-magnitude:unit}{unidades} próprias, ambas definidas a partir da própria escada.

\begin{definition}[Unidade astronômica e ano-luz]\label{def:g10:orders-of-magnitude:lightyear}
A \emph{unidade astronômica}\index{unidade astronômica}
($\qty{1}{au} = \qty{1.496e11}{m}$) é a distância média entre a Terra e o
Sol. O \emph{ano-luz}\index{ano-luz} (\unit{ly}) é a distância que a luz
percorre em um ano. A luz se propaga no vácuo a
\[
c = \qty{3.00e8}{m/s},
\]
a maior velocidade que a natureza permite --- de modo que o ano-luz é
uma distância, e não um tempo.
\end{definition}

\begin{omfigure}
\begin{tikzpicture}[x=1.08cm, font=\small]
  \draw[-Stealth] (7.6,0) -- (17.5,0);
  \foreach \e in {8,10,12,14,16} {
    \draw (\e,-0.12) -- (\e,0.12);
    \node[below, font=\footnotesize] at (\e,-0.14) {$10^{\e}$};
  }
  \foreach \x/\t/\d/\a in {8.58/Lua/{\qty{1.3}{s}}/north,
                           11.17/Sol/{\qty{8.3}{min}}/north,
                           12.65/Netuno/{\qty{4.2}{h}}/north,
                           15.98/{$\qty{1}{ly}$}/{\qty{1}{ano}}/north east,
                           16.60/Próxima/{\qty{4.2}{anos}}/north west} {
    \draw[omDef, thick] (\x,0.14) -- (\x,0.55);
    \node[rotate=52, anchor=south west, inner sep=1pt, omDef]
      at (\x,0.58) {\t};
    \node[anchor=\a, font=\footnotesize, omThm] at (\x,-0.55) {\d};
  }
  \node[right, font=\footnotesize] at (17.5,0) {distância (\unit{m})};
\end{tikzpicture}

{\small Distâncias a partir da Terra sobre a régua multiplicativa, com o
tempo de viagem da luz em vermelho embaixo de cada uma: a Lua está a
segundos-luz, os planetas a horas-luz, as estrelas a anos-luz.}
\end{omfigure}

\begin{example}[O ano-luz em metros]\label{ex:g10:orders-of-magnitude:ly}
Um ano tem $365.25 \times 24 \times 3600 = \qty{3.156e7}{s}$
(agradavelmente perto de $\pi \times 10^{7}\,\unit{s}$). Portanto
\[
\qty{1}{ly} = c \times \qty{3.156e7}{s}
= \qty{3.00e8}{m/s} \times \qty{3.156e7}{s}
= \qty{9.47e15}{m}.
\]
Próxima Centauri, a estrela mais próxima depois do Sol, está a
$\qty{4.2}{ly} \approx \qty{4.0e16}{m}$: a luz, que atravessa a
distância Terra--Sol em oito minutos, leva mais de quatro anos para
chegar de lá.
\end{example}

\begin{method}[Conferir pela ordem de grandeza]\label{met:g10:orders-of-magnitude:sanity}
Antes de confiar em qualquer resultado calculado:
\begin{enumerate}
  \item arredonde cada dado de entrada para a sua \omterm{def:g10:orders-of-magnitude:oom}{ordem de grandeza};
  \item refaça a conta só com potências de dez, usando a
        \cref{prop:g10:orders-of-magnitude:powers} --- isso leva
        segundos;
  \item compare: a resposta precisa tem de concordar com a estimativa a
        menos de um fator dez, mais ou menos, e a sua \omterm{def:g10:orders-of-magnitude:unit}{unidade} tem de ser
        a certa. Se qualquer uma das duas conferências falhar, cace o
        erro antes de seguir adiante.
\end{enumerate}
\end{method}

\begin{remark}[Por que se dar ao trabalho]\label{rem:g10:orders-of-magnitude:sanity}
Uma calculadora informa sem pestanejar que uma maçã tem massa
$\qty{4e8}{kg}$ --- ela não tem como saber. Nenhuma fórmula denuncia o
absurdo; o hábito denuncia. O físico estima primeiro e calcula depois: a
estimativa pega vírgulas fora do lugar, \omterm{def:g10:orders-of-magnitude:unit}{unidades} trocadas e frações
invertidas, os três erros por trás da maioria das respostas erradas em
qualquer nível.
\end{remark}

\section{Medida e incerteza}

Nenhum instrumento lê o valor verdadeiro. Uma medida só é honesta quando
informa também o quanto pode estar errada.

\begin{definition}[Incerteza absoluta e incerteza relativa]\label{def:g10:orders-of-magnitude:uncertainty}
Uma medida de uma grandeza $x$ se enuncia como
\[
x = x_{\text{m}} \pm \Delta x ,
\]
em que $x_{\text{m}}$ é o valor medido e a \emph{incerteza
absoluta}\index{incerteza!absoluta} $\Delta x > 0$, na mesma \omterm{def:g10:orders-of-magnitude:unit}{unidade},
limita o erro plausível: afirma-se que o valor verdadeiro está entre
$x_{\text{m}} - \Delta x$ e $x_{\text{m}} + \Delta x$. A \emph{incerteza
relativa}\index{incerteza!relativa} é a razão $\Delta x / x_{\text{m}}$,
em geral dada em porcentagem; é ela que a palavra ``preciso'' mede.
\end{definition}

\begin{method}[Ler um instrumento graduado]\label{met:g10:orders-of-magnitude:reading}
Para uma régua, uma proveta, um mostrador:
\begin{enumerate}
  \item leia a graduação mais próxima da marca, com o olho perpendicular
        à escala (um olhar de viés desloca a leitura);
  \item tome como \omterm{def:g10:orders-of-magnitude:uncertainty}{incerteza absoluta} pelo menos metade da menor
        graduação --- mais, se a borda do objeto ou a superfície do
        líquido estiver difusa;
  \item escreva o resultado como $x_{\text{m}} \pm \Delta x$ com a sua
        \omterm{def:g10:orders-of-magnitude:unit}{unidade}, e não guarde nenhum dígito além do duvidoso.
\end{enumerate}
\end{method}

\begin{omfigure}
\begin{tikzpicture}[x=0.9cm, font=\small]
  \fill[omThm, opacity=0.15] (6.5,-0.85) rectangle (7.5,1.45);
  \draw (-1.2,0) rectangle (10.6,1.1);
  \foreach \x in {0,...,10} \draw (\x,1.1) -- (\x,0.72);
  \foreach \x/\l in {0/120, 5/125, 10/130}
    \node[font=\footnotesize] at (\x,0.4) {\l};
  \fill[gray!25] (-1.0,-0.75) rectangle (7.0,-0.12);
  \draw (-1.0,-0.75) rectangle (7.0,-0.12);
  \draw[dashed, omThm] (7,-0.85) -- (7,1.45);
  \node[anchor=west, omThm] at (7.7,-0.5)
    {$L = (127.0 \pm 0.5)\,\unit{mm}$};
\end{tikzpicture}

{\small Leitura de uma régua milimetrada: a ponta da barra cai mais perto
da graduação de $\qty{127}{mm}$, e a faixa sombreada de meia graduação
de cada lado é a afirmação honesta.}
\end{omfigure}

\begin{example}[Precisão relativa]\label{ex:g10:orders-of-magnitude:reading}
A leitura acima tem \omterm{def:g10:orders-of-magnitude:uncertainty}{incerteza relativa}
$\frac{0.5}{127} \approx 0.4\%$. A mesma régua, medindo um parafuso de
$\qty{12}{mm}$, dá $\frac{0.5}{12} \approx 4\%$: o instrumento idêntico
é dez vezes menos preciso no objeto pequeno. A precisão é uma
propriedade da \emph{medida}, não da ferramenta.
\end{example}

\begin{remark}[Absoluta pequena, relativa enorme --- e vice-versa]\label{rem:g10:orders-of-magnitude:precision}
Estações a laser medem a distância Terra--Lua, $\qty{3.84e8}{m}$, com
cerca de $\pm\qty{3}{cm}$: uma \omterm{def:g10:orders-of-magnitude:uncertainty}{incerteza relativa} de $10^{-10}$, uma das
medidas mais finas já feitas --- e com um erro absoluto de que régua
nenhuma se gabaria. Pergunte sempre \emph{relativa a quê}: o par
$(\Delta x,\ \Delta x / x_{\text{m}})$ conta a história inteira, e cada
um dos dois números, sozinho, pode enganar.
\end{remark}

\section{Exercícios}

\begin{exercise}[$\star$]\label{exo:g10:orders-of-magnitude:1}
Converta para \omterm{def:g10:orders-of-magnitude:unit}{metros}, em \omterm{def:g10:orders-of-magnitude:scinot}{notação científica}:
\begin{enumerate}
  \item $\qty{750}{nm}$ (comprimento de onda da luz vermelha);
  \item $\qty{2.5}{km}$;
  \item $\qty{3.2}{\micro m}$ (uma bactéria do solo);
  \item $\qty{0.15}{Tm}$ (a distância Terra--Sol, num disfarce
        incomum).
\end{enumerate}
\end{exercise}

\begin{exercise}[$\star$]\label{exo:g10:orders-of-magnitude:2}
Escreva em \omterm{def:g10:orders-of-magnitude:scinot}{notação científica} e dê o número de \omterm{def:g10:orders-of-magnitude:sigfig}{algarismos
significativos}:
\begin{enumerate}
  \item a distância Terra--Sol, $\qty{149600000}{km}$, em \omterm{def:g10:orders-of-magnitude:unit}{metros};
  \item o diâmetro de um átomo de hidrogênio,
        $\qty{0.000000000106}{m}$;
  \item a população mundial, cerca de $\num{8100000000}$;
  \item uma duração de $\qty{0.00520}{s}$.
\end{enumerate}
\end{exercise}

\begin{exercise}[$\star$]\label{exo:g10:orders-of-magnitude:3}
Calcule sem calculadora, dando cada resultado em \omterm{def:g10:orders-of-magnitude:scinot}{notação científica}:
\[
(3\times10^{8})\times(2\times10^{-5}), \qquad
\frac{8\times10^{4}}{4\times10^{-3}}, \qquad
(5\times10^{6})^{2}, \qquad
\frac{(4\times10^{-3})\times(3\times10^{-9})}{6\times10^{-7}} .
\]
\end{exercise}

\begin{exercise}[$\star$]\label{exo:g10:orders-of-magnitude:4}
Dê a \omterm{def:g10:orders-of-magnitude:oom}{ordem de grandeza}, com a sua \omterm{def:g10:orders-of-magnitude:unit}{unidade}, de: uma altura humana de
$\qty{1.7}{m}$; o monte Everest, $\qty{8848}{m}$; uma bactéria
\emph{E.~coli}, $\qty{2e-6}{m}$; a massa da Terra, $\qty{5.97e24}{kg}$.
\end{exercise}

\begin{exercise}[$\star$]\label{exo:g10:orders-of-magnitude:5}
Uma régua é graduada em milímetros. Você mede um lápis e obtém
$\qty{127}{mm}$.
\begin{enumerate}
  \item Enuncie o resultado com a sua \omterm{def:g10:orders-of-magnitude:uncertainty}{incerteza absoluta}.
  \item Calcule a \omterm{def:g10:orders-of-magnitude:uncertainty}{incerteza relativa}.
  \item A mesma régua mede uma borracha de $\qty{12}{mm}$. Qual das duas
        medidas é a mais precisa, e por que fator?
\end{enumerate}
\end{exercise}

\begin{exercise}[$\star\star$]\label{exo:g10:orders-of-magnitude:6}
Usando $c = \qty{3.00e8}{m/s}$ e $\qty{1}{ly} = \qty{9.47e15}{m}$:
\begin{enumerate}
  \item Exprima em \omterm{def:g10:orders-of-magnitude:unit}{metros} as distâncias até Próxima Centauri
        ($\qty{4.24}{ly}$) e até Sirius ($\qty{8.6}{ly}$).
  \item A sonda \emph{Voyager 1} está a cerca de $\qty{2.4e13}{m}$ da
        Terra. Quanto tempo leva o seu sinal de rádio, que viaja a $c$,
        para chegar até nós? Exprima a resposta em horas.
\end{enumerate}
\end{exercise}

\begin{exercise}[$\star\star$]\label{exo:g10:orders-of-magnitude:7}
Quanto tempo a luz leva para chegar à Terra vinda da Lua
($\qty{3.84e8}{m}$), do Sol ($\qty{1.496e11}{m}$) e de Netuno
($\qty{4.50e12}{m}$)? Dê cada resposta em uma \omterm{def:g10:orders-of-magnitude:unit}{unidade} sensata e complete
a frase: ``Quando você olha para Netuno, vê-o como ele era\dots''
\end{exercise}

\begin{exercise}[$\star\star$]\label{exo:g10:orders-of-magnitude:8}
Uma resma de $500$ folhas de papel, medida com uma régua milimetrada,
tem $\qty{52.0}{mm}$ de espessura, com incerteza $\pm\qty{0.5}{mm}$.
\begin{enumerate}
  \item Deduza a espessura de uma folha, com a sua \omterm{def:g10:orders-of-magnitude:uncertainty}{incerteza absoluta}.
  \item Qual é a \omterm{def:g10:orders-of-magnitude:uncertainty}{incerteza relativa} dessa espessura?
  \item Por que medir a resma é enormemente melhor do que medir uma
        folha diretamente com a mesma régua? Quantifique.
\end{enumerate}
\end{exercise}

\begin{exercise}[$\star\star$]\label{exo:g10:orders-of-magnitude:9}
Tome $\qty{2e-10}{m}$ como diâmetro de um átomo. Quantos átomos, lado a
lado, cobrem: uma bactéria de $\qty{2.0}{\micro m}$; um fio de cabelo de
$\qty{70}{\micro m}$; um ser humano de $\qty{1.7}{m}$? Dê também a
última resposta como \omterm{def:g10:orders-of-magnitude:oom}{ordem de grandeza}.
\end{exercise}

\begin{exercise}[$\star\star$]\label{exo:g10:orders-of-magnitude:10}
Uma maquete representa o Sol (diâmetro $\qty{1.39e9}{m}$) por uma bola
de futebol de diâmetro $\qty{22}{cm}$.
\begin{enumerate}
  \item Calcule o fator de escala da maquete.
  \item Encontre o diâmetro da Terra na maquete ($\qty{1.27e7}{m}$) e a
        sua distância à bola (distância real $\qty{1.496e11}{m}$).
  \item A que distância da Terra da maquete fica a Lua da maquete
        (distância real $\qty{3.84e8}{m}$)? E a que distância fica
        Netuno da maquete (distância real $\qty{4.50e12}{m}$)?
\end{enumerate}
\end{exercise}

\begin{exercise}[$\star\star$]\label{exo:g10:orders-of-magnitude:11}
\omterm{def:g10:orders-of-magnitude:sigfig}{Algarismos significativos} em ação.
\begin{enumerate}
  \item Uma folha A4 mede $\qty{21.0}{cm}$ por $\qty{29.7}{cm}$. Dê a
        sua área com o número correto de \omterm{def:g10:orders-of-magnitude:sigfig}{algarismos significativos}.
  \item Um velocista percorre $\qty{100.0}{m}$ em $\qty{12.3}{s}$. Dê a
        velocidade média corretamente arredondada.
  \item Explique em uma frase o que haveria de desonesto em anunciar a
        área como $\qty{623.70}{cm^2}$.
\end{enumerate}
\end{exercise}

\begin{exercise}[$\star\star\star$]\label{exo:g10:orders-of-magnitude:12}
Uma gota de óleo de volume $\qty{1.0}{mm^3}$ é depositada sobre água
parada e se espalha num filme circular de diâmetro $\qty{1.1}{m}$ --- um
filme que capítulos posteriores mostrarão ter a espessura de uma única
molécula.
\begin{enumerate}
  \item Calcule a área do filme e, em seguida, a sua espessura.
  \item Deduza uma \omterm{def:g10:orders-of-magnitude:oom}{ordem de grandeza} para o tamanho de uma molécula e
        confira-a com o degrau do átomo na escada.
\end{enumerate}
(Essa experiência de fundo de envelope, tentada pela primeira vez por
Benjamin Franklin num lago de Londres, foi uma das primeiras medidas
honestas que a humanidade fez do mundo molecular.)
\end{exercise}

\begin{exercise}[$\star\star\star$]\label{exo:g10:orders-of-magnitude:13}
Estime o número de batimentos cardíacos de uma vida humana. Enuncie as
suas hipóteses (frequência cardíaca em repouso, duração da vida),
mantenha a aritmética em potências de dez e um dígito de \omterm{def:g10:orders-of-magnitude:scinot}{mantissa}, e dê
a resposta final apenas como \omterm{def:g10:orders-of-magnitude:oom}{ordem de grandeza}. Por que dá-la com quatro
\omterm{def:g10:orders-of-magnitude:sigfig}{algarismos significativos} não faria sentido?
\end{exercise}

\begin{exercise}[$\star\star\star$]\label{exo:g10:orders-of-magnitude:14}
Uma folha de papel tem $\qty{0.10}{mm}$ de espessura, e cada dobra
duplica a espessura da pilha. Usando $2^{10} \approx 10^{3}$:
\begin{enumerate}
  \item Depois de quantas dobras a pilha ultrapassaria o monte Everest
        ($\qty{8848}{m}$)?
  \item Depois de quantas dobras ela alcançaria a Lua
        ($\qty{3.84e8}{m}$)?
  \item Na prática, uma folha não se dobra mais do que umas sete vezes.
        Isso invalida a aritmética ou apenas a experiência?
\end{enumerate}
\end{exercise}

\begin{exercise}[$\star\star\star$]\label{exo:g10:orders-of-magnitude:15}
A telemetria a laser dá a distância Terra--Lua como
$\qty{3.84e8}{m} \pm \qty{3}{cm}$; uma boa régua dá uma mesa como
$\qty{1.000}{m} \pm \qty{1}{mm}$.
\begin{enumerate}
  \item Calcule as duas \omterm{def:g10:orders-of-magnitude:uncertainty}{incertezas relativas}. Qual medida é a mais
        precisa, e por que fator?
  \item Com que \omterm{def:g10:orders-of-magnitude:uncertainty}{incerteza absoluta} seria preciso medir a mesa para
        igualar a precisão relativa da medida da Lua? Compare esse
        comprimento com o tamanho de um átomo e conclua com o devido
        respeito.
\end{enumerate}
\end{exercise}

\section{Problema: do átomo a Andrômeda}

\begin{problem}[{Problema de fim de semana --- um único grão de areia
dimensiona o átomo, o Sistema Solar, a galáxia e o universo observável, e
localiza você na escada entre eles}]
\label{pb:g10:orders-of-magnitude:1}
Um grão de areia é o objeto mais humilde da física, e neste fim de
semana ele faz hora extra: primeiro contamos os seus átomos, depois
encolhemos o Sol até o tamanho dele e medimos o cosmos a passos, em
seguida deixamos a luz fazer a topografia e, por fim, contamos os átomos
de \emph{você}. Dados, usados o tempo todo: diâmetro do grão
$\qty{0.5}{mm}$; diâmetro do átomo $\qty{2.5e-10}{m}$; densidade da
areia $\qty{2500}{kg/m^3}$; diâmetros: Sol $\qty{1.39e9}{m}$, Terra
$\qty{1.27e7}{m}$, Via Láctea $\qty{9.5e20}{m}$, universo observável
$\qty{8.8e26}{m}$; distâncias até nós: Lua $\qty{3.84e8}{m}$, Sol
$\qty{1.496e11}{m}$, Netuno $\qty{4.5e12}{m}$, Próxima Centauri
$\qty{4.0e16}{m}$, galáxia de Andrômeda $\qty{2.4e22}{m}$;
$c = \qty{3.00e8}{m/s}$; uma molécula de água tem massa
$\qty{3.0e-26}{kg}$ e contém três átomos.

\medskip
\noindent\textbf{Parte I --- Anatomia do grão.}
\begin{enumerate}
  \item Escreva o diâmetro do grão e o diâmetro do átomo em \omterm{def:g10:orders-of-magnitude:unit}{metros}, em
        \omterm{def:g10:orders-of-magnitude:scinot}{notação científica}, e diga de qual \omterm{def:g10:orders-of-magnitude:prefixes}{prefixo} do SI cada um está
        mais perto.
  \item Quantos átomos cabem lado a lado ao longo de um diâmetro do
        grão?
  \item Modele o grão como um cubo de aresta $\qty{0.5}{mm}$ empacotado
        com átomos. Estime o número de átomos do grão, como \omterm{def:g10:orders-of-magnitude:oom}{ordem de
        grandeza}.
  \item Calcule o volume desse cubo e, em seguida, a massa do grão.
        Exprima a massa em miligramas.
  \item Confira o modelo: divida a massa do grão pelo seu número de
        átomos e compare com a \omterm{def:g10:orders-of-magnitude:oom}{ordem de grandeza} conhecida das massas
        atômicas (alguns $10^{-26}\,\unit{kg}$). Veredito?
\end{enumerate}

\medskip
\noindent\textbf{Parte II --- O Sol como grão de areia.}
Construímos uma maquete do cosmos em que o Sol é o grão: cada tamanho da
maquete é o tamanho real multiplicado por um fator $s$.
\begin{enumerate}[resume]
  \item Calcule o fator de escala $s$.
  \item Na maquete, encontre o diâmetro da Terra e a sua distância ao
        grão-Sol. Que objeto do mundo da Parte I tem, mais ou menos, o
        tamanho da Terra da maquete?
  \item A que distância do grão orbita o Netuno da maquete? Que móvel
        comportaria o sistema planetário inteiro?
  \item A que distância está a Próxima Centauri da maquete?
  \item Dois grãos de areia, a $\qty{14}{km}$ um do outro, e
        praticamente nada no meio: diga em uma ou duas frases o que a
        maquete revela sobre o quanto a galáxia é vazia.
\end{enumerate}

\medskip
\noindent\textbf{Parte III --- A luz faz a topografia.}
\begin{enumerate}[resume]
  \item Quanto tempo a luz leva para atravessar o grão de areia real?
  \item Quanto tempo a luz leva para chegar até nós vinda da Lua e vinda
        do Sol?
  \item Recalcule o \omterm{def:g10:orders-of-magnitude:lightyear}{ano-luz} em \omterm{def:g10:orders-of-magnitude:unit}{metros} a partir de $c$ e da duração do
        ano, e verifique que a distância até Próxima que consta dos
        dados corresponde a cerca de $\qty{4.2}{ly}$.
  \item Converta a distância até Andrômeda em anos-luz. Aproximadamente
        quando partiu de Andrômeda a luz que agora entra no seu olho, e
        o que andava sobre a Terra naquele momento?
  \item Que distância a luz percorre em um nanossegundo? Um processador
        com relógio de $\qty{3}{GHz}$ inicia uma operação a cada terço
        de nanossegundo: explique em uma frase por que um computador
        rápido tem de ser um computador pequeno.
\end{enumerate}

\medskip
\noindent\textbf{Parte IV --- Você, entre o átomo e Andrômeda.}
\begin{enumerate}[resume]
  \item Um ser humano é, grosso modo, $\qty{70}{kg}$ de água. Usando os
        dados, estime a massa média de um átomo na água e, em seguida, o
        número de átomos de um ser humano.
  \item Compare essa contagem com a dos átomos do grão (questão 3) e com
        as cerca de $2\times10^{23}$ estrelas do universo observável. Dê
        os dois fatores.
  \item Na escada, você está mais perto do átomo ou do universo
        observável? Compare as razões $\text{(você)}/\text{(átomo)}$ e
        $\text{(universo)}/\text{(você)}$, tomando $\qty{1.7}{m}$ para
        você.
  \item Encontre o tamanho $x$ que fica exatamente no meio da escada,
        entre o núcleo ($\qty{e-15}{m}$) e o universo observável: $x$
        tem de satisfazer $x / 10^{-15} = 10^{26.9} / x$ (usando
        $\qty{8.8e26}{m} \approx 10^{26.9}\,\unit{m}$). Que objeto do
        Sistema Solar tem, por acaso, justamente esse tamanho?
  \item O desfecho. Encolha o universo observável pelo fator $s$ da
        Parte II, de modo que o Sol seja um grão de areia. Dê o diâmetro
        do universo da maquete em \omterm{def:g10:orders-of-magnitude:unit}{metros} e depois em \omterm{def:g10:orders-of-magnitude:lightyear}{unidades
        astronômicas}, e conclua em uma frase: mesmo com cada estrela
        virando um grão de areia, o universo transborda de toda escala
        humana --- só as potências de dez o seguram.
\end{enumerate}
\end{problem}
